\chapter{Erzeugte Äquivalenzen, Vereinigung und Verkettung}
\label{b11guide:ch:generated}

\section{Die von Paaren erzeugte Äquivalenz}

\begin{definition}[Erzeugte Äquivalenzrelation]
\label{b11guide:def:generated}
Für $R\subseteq A\times A$ sei
\[
 \EqHull_A(R)=
 \bigcap\{E\in\EqSet(A)\mid R\subseteq E\}.
\]
Die Familie unter dem Durchschnitt ist nichtleer, denn sie
enthält $\nabla_A$.
$\EqHull_A(R)$ heißt die \emph{von $R$ erzeugte
Äquivalenzrelation} oder die \emph{Äquivalenzhülle von $R$}.
Die Ausgangsmenge $R$ darf beliebige vorgegebene Paare enthalten;
Gegenstand der folgenden Sätze ist ihre Äquivalenzhülle.
\end{definition}

\begin{theorem}[Existenz, Minimalität und Erkennung]
\label{b11guide:thm:hull}
Für $R\subseteq A\times A$ gilt
\[
 \EqHull_A(R)\in\EqSet(A),\qquad R\subseteq\EqHull_A(R).
\]
Für jedes $E\in\EqSet(A)$ gilt
\[
 R\subseteq E\quad\Longleftrightarrow\quad\EqHull_A(R)\subseteq E.
\]
Insbesondere ist die Hülle die eindeutig bestimmte kleinste
Äquivalenzrelation, die $R$ enthält, und
\[
 R\in\EqSet(A)\quad\Longleftrightarrow\quad\EqHull_A(R)=R.
\]
\end{theorem}

\begin{theorem}[Beschreibung durch endliche Verbindungsketten]
\label{b11guide:thm:chains}
Für $R\subseteq A\times A$ und $x,y\in A$ gilt
$x\sim_{\EqHull_A(R)}y$ genau dann, wenn es eine endliche Kette
\[
 x=x_0,x_1,\ldots,x_n=y\qquad(n\geq0)
\]
in $A$ gibt, bei der für jedes $i<n$ gilt
\[
 (x_i,x_{i+1})\in R\quad\text{oder}\quad(x_{i+1},x_i)\in R.
\]
Die Kette der Länge $0$ ist zugelassen und bedeutet $x=y$.
Die Klasse $[x]_{\EqHull_A(R)}$ besteht genau aus den so von $x$
aus erreichbaren Elementen.
\end{theorem}

\begin{theorem}[Gesetze der Äquivalenzhülle]
Für $R,S\subseteq A\times A$ gilt
\[
\begin{aligned}
 R\subseteq S&\ \Rightarrow\ \EqHull_A(R)\subseteq\EqHull_A(S),\\
 \EqHull_A(\EqHull_A(R))&=\EqHull_A(R),\\
 \EqHull_A(\varnothing)&=\Delta_A,\\
 \EqHull_A(R\cup\Delta_A)&=\EqHull_A(R),\\
 \EqHull_A(R\cup S)
   &=\EqHull_A\bigl(\EqHull_A(R)\cup\EqHull_A(S)\bigr).
\end{aligned}
\]
\end{theorem}

\begin{theorem}[Gesättigtheit anhand der erzeugenden Paare]
Seien $R\subseteq A\times A$ und $B\subseteq A$. Dann gilt
\[
\begin{gathered}
 B\text{ ist }\EqHull_A(R)\text{-gesättigt}\\
 \Longleftrightarrow\
 \forall(x,y)\in R\;(x\in B\Longleftrightarrow y\in B).
\end{gathered}
\]
\end{theorem}

\section{Vereinigungen von Äquivalenzrelationen}

\begin{theorem}[Gemeinsam erzeugte Äquivalenzrelation]
\label{b11guide:thm:generatedfamily}
Für $\mathcal E\subseteq\EqSet(A)$ ist
\[
 H=\EqHull_A\left(\bigcup\mathcal E\right)
\]
die kleinste Äquivalenzrelation auf $A$, die jedes $E\in\mathcal E$
enthält. Für $x,y\in A$ gilt $x\sim_H y$ genau dann, wenn es
eine endliche Kette von $x$ nach $y$ gibt, deren benachbarte
Elemente jeweils in einer Relation aus $\mathcal E$ äquivalent
sind. Die benutzte Relation darf von Schritt zu Schritt wechseln.
Für $\mathcal E=\varnothing$ gilt $H=\Delta_A$.
\end{theorem}

\begin{theorem}[Genaues Vereinigungskriterium]
\label{b11guide:thm:union}
Für $E,F\in\EqSet(A)$ sind äquivalent:
\[
\begin{aligned}
 &E\cup F\in\EqSet(A),\\
 &E\cup F\text{ ist transitiv},\\
 &\forall C\in A/E\;\forall D\in A/F:\\[-.4em]
 &\hspace{2em}
 C\cap D\ne\varnothing\Rightarrow(C\subseteq D\ \lor\ D\subseteq C).
\end{aligned}
\]
Insbesondere genügt $E\subseteq F$ oder $F\subseteq E$.
Ohne eine solche Bedingung kann die Vereinigung die
Transitivität verlieren.
\end{theorem}

\begin{theorem}[Vereinigung mit gemeinsamen umfassenden Familiengliedern]
Sei $\varnothing\ne\mathcal E\subseteq\EqSet(A)$. Falls
\[
 \forall E,F\in\mathcal E\;\exists G\in\mathcal E\;
                                      (E\cup F\subseteq G),
\]
so gilt
\[
 \bigcup\mathcal E\in\EqSet(A).
\]
\end{theorem}

\section{Wann ist eine Verkettung wieder eine Äquivalenz?}

\begin{definition}[Umkehrrelation und Verkettung]
Für $R,S\subseteq A\times A$ seien
\[
\begin{aligned}
 R^{-1}&=\{(y,x)\mid(x,y)\in R\},\\
 S\circ R&=\{(x,z)\in A\times A\mid
               \exists y\in A\;((x,y)\in R\land(y,z)\in S)\}.
\end{aligned}
\]
Bei $S\circ R$ wird zuerst $R$, dann $S$ durchlaufen.
Die Definitionen werden nur für die folgenden
Äquivalenzkriterien verwendet.
\end{definition}

\begin{theorem}[Charakterisierung durch Mengenoperationen]
Für $R\subseteq A\times A$ gilt
\[
 R\in\EqSet(A)
 \quad\Longleftrightarrow\quad
 \Delta_A\subseteq R,\quad R^{-1}=R,\quad R\circ R=R.
\]
\end{theorem}

\begin{theorem}[Vertauschbare Äquivalenzrelationen]
\label{b11guide:thm:permuting}
Für $E,F\in\EqSet(A)$ sind äquivalent:
\[
\begin{aligned}
 &E\circ F=F\circ E,\\
 &F\circ E\text{ ist symmetrisch},\\
 &F\circ E\text{ ist transitiv},\\
 &F\circ E\in\EqSet(A).
\end{aligned}
\]
Unter diesen Bedingungen gilt
\[
 F\circ E=E\circ F=\EqHull_A(E\cup F).
\]
Für $x\in A$ gilt dann außerdem
\[
 [x]_{F\circ E}=\Sat_F([x]_E)=\Sat_E([x]_F).
\]
\end{theorem}

